\documentclass{xourse}


\input{../../preamble.tex}


\definecolor{fvdnavy}{HTML}{071D43}
\definecolor{fvdcyan}{HTML}{20BFC6}
\definecolor{fvdcyanlight}{HTML}{EFFCFB}
\definecolor{fvdmagenta}{HTML}{F10D69}
\definecolor{fvdmagentalight}{HTML}{FFF1F7}
\definecolor{fvdtext}{HTML}{163044}





%=================================================
% Pedagogische omgevingen
% PDF: tcolorbox
% HTML: learning-box
%=================================================

\newenvironment{voorkennis}
{%
    \ifdefined\HCode
        \HCode{%
<section class="learning-box learning-box--prior">
<div class="learning-box__header">
<span class="learning-box__icon" aria-hidden="true"></span>
<span class="learning-box__title">Voorkennis</span>
</div>
<div class="learning-box__content">
        }%
        \begin{itemize}
    \else
       \begin{tcolorbox}[
    enhanced,
    breakable,
    colback=fvdcyanlight,
    colframe=fvdcyan,
    coltext=fvdtext,
    boxrule=0.5pt,
    leftrule=4pt,
    arc=4mm,
    outer arc=4mm,
    left=7mm,
    right=6mm,
    top=5mm,
    bottom=5mm,
    before skip=6mm,
    after skip=6mm,
    title=Voorkennis,
    coltitle=fvdcyan,
    fonttitle=\bfseries\Large,
    attach title to upper={\par\smallskip},
    boxed title style={
        empty,
        left=0mm,
        right=0mm,
        top=0mm,
        bottom=0mm
    },
    overlay={
        \draw[
            fvdcyan,
            line width=0.5pt
        ]
        ([xshift=42mm,yshift=-13mm]frame.north west)
        --
        ([xshift=-7mm,yshift=-13mm]frame.north east);
    }
]
        \begin{itemize}
    \fi
}
{%
    \end{itemize}
    \ifdefined\HCode
        \HCode{%
</div>
</section>
        }%
    \else
        \end{tcolorbox}
    \fi
}


\newenvironment{leerdoelen}
{%
    \ifdefined\HCode
        \HCode{%
<section class="learning-box learning-box--goals">
<div class="learning-box__header">
<span class="learning-box__icon" aria-hidden="true"></span>
<span class="learning-box__title">Leerdoelen</span>
</div>
<div class="learning-box__content">
        }%
        \begin{itemize}
    \else
\begin{tcolorbox}[
    enhanced,
    breakable,
    colback=fvdmagentalight,
    colframe=fvdmagenta,
    coltext=fvdtext,
    boxrule=0.5pt,
    leftrule=4pt,
    arc=4mm,
    outer arc=4mm,
    left=7mm,
    right=6mm,
    top=5mm,
    bottom=5mm,
    before skip=6mm,
    after skip=6mm,
    title=Leerdoelen,
    coltitle=fvdmagenta,
    fonttitle=\bfseries\Large,
    attach title to upper={\par\smallskip},
    boxed title style={
        empty,
        left=0mm,
        right=0mm,
        top=0mm,
        bottom=0mm
    },
    overlay={
        \draw[
            fvdmagenta,
            line width=0.5pt
        ]
        ([xshift=42mm,yshift=-13mm]frame.north west)
        --
        ([xshift=-7mm,yshift=-13mm]frame.north east);
    }
]
        \begin{itemize}
    \fi
}
{%
    \end{itemize}
    \ifdefined\HCode
        \HCode{%
</div>
</section>
        }%
    \else
        \end{tcolorbox}
    \fi
}


\addPrintStyle{../..}
  

\fvdtitle{De digitale wereld begrijpen}
\author{Ingmar Herreman}


\begin{abstract}

    In deze module bouwen we de wiskundige basis op die je nodig hebt voor de
    rest van de cursus. We vertrekken van getallen en rekenvaardigheid en groeien
    stap voor stap naar algebra, formules en eenvoudige wiskundige modellen.
    
    \\ \ \\
    
    We oefenen hoofdrekenen, negatieve getallen, kommagetallen, breuken,
    procenten en verhoudingen. Daarna onderzoeken we machten, wortels en
    wetenschappelijke notatie. Vanuit concrete berekeningen maken we vervolgens
    de stap naar rekenen met letters en formules.
    \\ \ \\
    De wiskunde krijgt betekenis in toepassingen met onder andere pixels,
    resoluties, beeldverhoudingen, coördinaten, schaalfactoren en andere
    grootheden uit digitale creatie en technologie. Waar het zinvol is,
    gebruiken we GeoGebra om verbanden te onderzoeken en zichtbaar te maken.
    \\ \ \\
    Een computer kan razendsnel rekenen, maar jij moet kunnen beoordelen of een
    resultaat logisch is. Daarom leren we ook schatten, controleren, een geschikte
    strategie kiezen en kritisch naar resultaten kijken.
    
    \begin{center}
    \textbf{Het doel is niet alleen correct rekenen, maar wiskunde gebruiken om
    verbanden te begrijpen, problemen op te lossen en de digitale wereld te
    beschrijven.}
    \end{center}
      
  \end{abstract}


% ============================================================
% MODULEGEGEVENS
% ============================================================

\def\fvdmodulenummer{1}
\def\fvdmoduletitel{De digitale wereld begrijpen}
\def\fvdmodulebeeld{../images/banners/digitale_wereld.png}
\def\fvdschoolembleem{../images/banners/logo_athena.png}
\def\fvdschoollogo{../images/banners/logo_athena.png}
\renewcommand{\fvdschoolnaam}{Athena Pottelberg Kortrijk}
\def\vaknaam{Wiskunde}

\def\fvdklas{7 FVD}
\def\fvdschooljaar{2026--2027}
%\def\fvdmodulevariant{                             }
%\def\fvdmodulesubtitel{}



% ============================================================
% THEORIE + BIJHORENDE OEFENINGEN
%
% HTML:
%   theorie + oefeningen
%
% PDF:
%   alleen theorie
% ============================================================

\newcommand{\FVDchapterpair}[2]{%
  \ifdefined\HCode

    % ONLINE: theorie
    \activity{#1}%

    % ONLINE: oefeningen, indien aanwezig
    \if\relax\detokenize{#2}\relax
    \else
      \activity{#2}%
    \fi

  \else

    % PDF: alleen theorie
    \activity{#1}%

  \fi
}


\begin{document}

\FVDmodulecover



%%%%%%%%%%%%%%%%%%%%%%%%%%%%%%%%%%%%%%%%%%%%%%%%%%%%%%%%%%%%
%% MODULE-OPENING
%%%%%%%%%%%%%%%%%%%%%%%%%%%%%%%%%%%%%%%%%%%%%%%%%%%%%%%%%%%%

\ifdefined\HCode
  % Website: gebruik de gewone Ximera-titelpagina.
  \maketitle
\else
  % PDF: gebruik de automatisch gegenereerde module-opening.
  \IfFileExists{module-opening.generated.tex}{%
    % ============================================================
% AUTOMATISCH GEGENEREERD BESTAND
% Niet handmatig aanpassen.
% ============================================================

\ifdefined\HCode
  % Geen afzonderlijke module-openingspagina in HTML.
\else
  \FVDmoduleopening
    {De digitale wereld begrijpen}
    {In deze module bouwen we de wiskundige basis op die je nodig hebt voor de rest van de cursus. We vertrekken van getallen en rekenvaardigheid en groeien stap voor stap naar algebra, formules en eenvoudige wiskundige modellen. \\ \ \\ We oefenen hoofdrekenen, negatieve getallen, kommagetallen, breuken, procenten en verhoudingen. Daarna onderzoeken we machten, wortels en wetenschappelijke notatie. Vanuit concrete berekeningen maken we vervolgens de stap naar rekenen met letters en formules. \\ \ \\ De wiskunde krijgt betekenis in toepassingen met onder andere pixels, resoluties, beeldverhoudingen, coördinaten, schaalfactoren en andere grootheden uit digitale creatie en technologie. Waar het zinvol is, gebruiken we GeoGebra om verbanden te onderzoeken en zichtbaar te maken. \\ \ \\ Een computer kan razendsnel rekenen, maar jij moet kunnen beoordelen of een resultaat logisch is. Daarom leren we ook schatten, controleren, een geschikte strategie kiezen en kritisch naar resultaten kijken. \begin{center} \textbf{Het doel is niet alleen correct rekenen, maar wiskunde gebruiken om verbanden te begrijpen, problemen op te lossen en de digitale wereld te beschrijven.} \end{center}}
    {%
    \FVDmodulethemetile{1}{Rekenen in een digitale wereld}
    \FVDmodulethemetile{2}{Machten, wortels en grootteorde}
    \FVDmodulethemetile{3}{Algebra als taal}
    \FVDmodulethemetile{4}{Wiskunde onderzoeken en toepassen}
    }
\fi
%
  }{%
    \PackageWarning{FVD}{Automatische module-opening ontbreekt}%
    \maketitle
  }%
\fi


\FVDpart{Rekenen in een digitale wereld}

\FVDthemeopening{../images/banners/thema1-rekenen.png}



















% BEGIN AUTO THEME OVERVIEW
\begin{hoofdstukoverzicht}
  \FVDthemechaptertile{1}{Getallen en hoofdrekenen}
  \FVDthemechaptertile{2}{2 Breuken}
  \FVDthemechaptertile{3}{3 Procenten}
  \FVDthemechaptertile{4}{4 Verhoudingen}
\end{hoofdstukoverzicht}
% END AUTO THEME OVERVIEW

\begin{themaintro}{\textbf{Rekenvaardigheden zijn niet overbodig}}

  Computers voeren elke seconde enorme aantallen berekeningen uit. Games,
  animaties, beelden en digitale modellen bestaan uiteindelijk uit getallen:
  pixelafmetingen, kleurwaarden, coördinaten, tijdstappen, schaalfactoren,
  snelheden en nog veel meer.
  \\ \ \\
  Maar een computer die kan rekenen, maakt rekenvaardigheid niet overbodig.
  Integendeel. Wie met digitale technologie werkt, moet kunnen inschatten of
  een resultaat logisch is, fouten herkennen en begrijpen welke berekening
  achter een toepassing zit.
  \\ \ \\
  In dit thema bouwen we daarom een stevige rekenbasis op. We starten met
  getallen en hoofdrekenen en ontwikkelen strategieën om berekeningen handig,
  vlot en betrouwbaar uit te voeren. We leren schatten en controleren, zodat
  we niet alleen een antwoord kunnen berekenen, maar ook kunnen beoordelen of
  dat antwoord redelijk is.
  \\ \ \\
  Daarna onderzoeken we breuken, procenten en verhoudingen. Ze laten ons
  dezelfde hoeveelheid op verschillende manieren beschrijven en zijn
  onmisbaar bij onder andere schaalfactoren, beeldverhoudingen, resoluties en
  relatieve veranderingen.
  \\ \ \\
  Ook machten, wortels en wetenschappelijke notatie krijgen een plaats.
  Daarmee kunnen we zeer grote en zeer kleine getallen efficiënt voorstellen,
  vergelijken en verwerken.
  \\ \ \\
  Geleidelijk maken we vervolgens de overgang van rekenen met concrete
  getallen naar rekenen met letters en formules. Algebra maakt het mogelijk
  om patronen en verbanden algemeen te beschrijven en vormt daarmee een brug
  naar de wiskunde die later in de cursus aan bod komt.
  \\ \ \\
  Doorheen het thema verbinden we deze technieken met situaties uit de
  digitale wereld. We rekenen met pixels, afmetingen, coördinaten,
  schaalfactoren en andere grootheden die voorkomen in digitale creatie en
  technologie.
  \\ \ \\
  Het doel is daarbij niet om zoveel mogelijk berekeningen uit het hoofd te
  maken. We willen leren wanneer een berekening eenvoudig kan, welke strategie
  efficiënt is, wanneer een digitaal hulpmiddel nuttig is en hoe we het
  verkregen resultaat kritisch kunnen controleren.
  \\ \ \\
  Zo groeien getallen uit van iets waarmee we rekenen tot een taal waarmee we
  de digitale wereld kunnen beschrijven en begrijpen.

\end{themaintro}










\FVDchapterpair{Thema1-Rekenen/1-Getallen}{}
\FVDchapterpair{Thema1-Rekenen/2-Breuken}{}
\FVDchapterpair{Thema1-Rekenen/3-Procenten}{}
\FVDchapterpair{Thema1-Rekenen/4-Verhoudingen}{}



%%%%%%%%%%%%%%%%%%%%%%%%%%%%%%%%%%%%%%%%%%%%%%%%%%%%%%%%%%%%
%% THEMA 2
%%%%%%%%%%%%%%%%%%%%%%%%%%%%%%%%%%%%%%%%%%%%%%%%%%%%%%%%%%%

\FVDpart{Machten, wortels en grootteorde}

















% BEGIN AUTO THEME OVERVIEW
\begin{hoofdstukoverzicht}
  \FVDthemeoverviewempty
\end{hoofdstukoverzicht}
% END AUTO THEME OVERVIEW

\FVDpart{Algebra als taal}
















% BEGIN AUTO THEME OVERVIEW
\begin{hoofdstukoverzicht}
  \FVDthemeoverviewempty
\end{hoofdstukoverzicht}
% END AUTO THEME OVERVIEW

\FVDpart{Wiskunde onderzoeken en toepassen}




































\end{document}


% BEGIN AUTO THEME OVERVIEW
\begin{hoofdstukoverzicht}
  \FVDthemeoverviewempty
\end{hoofdstukoverzicht}
% END AUTO THEME OVERVIEW
